% --- Archon physics formalization source begin ---
% archon:physics
% archon:covers PhyXMiniProblems/problem_phyx_mini_0944.lean
% archon:source-report reports/phyx_mini/problem_phyx_mini_0944.source.json

\chapter{Physics problem phyx\_mini\_0944}
\label{ch:PhyXMiniProblems_problem_phyx_mini_0944}

\paragraph{Problem source.}
\#\# Physical scenario

A radio station on the earth’s surface emits a sinusoidal wave with average total power 50\textbackslash{} \textbackslash{}mathrm\{kW\} (figure). Assuming that the transmitter radiates equally in all directions above the ground (which is unlikely in real situations)

\#\# Question

find the electric-field and magnetic-field amplitudes \textbackslash{}( E\_\textbackslash{}mathrm\{max\} \textbackslash{}) and \textbackslash{}( B\_\textbackslash{}mathrm\{max\} \textbackslash{}) detected by a satellite 100\textbackslash{} \textbackslash{}mathrm\{km\} from the antenna.

\#\# Answer choices

- A: A: 2.2 \textbackslash{}times 10\textasciicircum{}\{-11\}T

- B: B: 8.17 \textbackslash{}times 10\textasciicircum{}\{-11\}T

- C: C: 1 \textbackslash{}times 10\textasciicircum{}\{-11\}T

- D: D: 2.6 \textbackslash{}times 10\textasciicircum{}\{-11\}T

\#\# Figure caption (auxiliary; use the image as primary evidence)

The image depicts a diagram showing a communication setup between a transmitter and a satellite. The scene includes:

- **Transmitter**: Located at the base of a semi-transparent blue dome, indicating a coverage area.
- **Satellite**: Positioned above the dome, shown as an orbiting object with solar panels.
- **Dome**: Represents the coverage area with a radius noted as “\textbackslash{}( r = 100 \textbackslash{}, \textbackslash{}text\{km\} \textbackslash{}).”
- **Arrow**: A black arrow connects the transmitter and the satellite, indicating the direction of the signal.
- **Text**: Labels “Transmitter” and “Satellite” identify the respective objects.

Overall, the image illustrates a communication link and coverage area between a transmitter and a satellite.

\#\# Dataset metadata

Electromagnetic Waves; Physical Model Grounding Reasoning, Multi-Formula Reasoning

\paragraph{Recorded answer/context.}
B: B: 8.17 \textbackslash{}times 10\textasciicircum{}\{-11\}T

\paragraph{Figure/image path.}
/root/proposal\_for\_physic/science-mango/phyx\_mini\_run/phyx\_data/test\_image/944.png

\paragraph{Formalization target.}
create a compiling Lean file with sorry bodies at `PhyXMiniProblems/problem\_phyx\_mini\_0944.lean`. The Lean declarations must preserve the physical quantities, dimensions or dimensional roles, figure labels, governing-law hypotheses, and final relation expressed by this problem.
Use Mathlib/Physlib names found through LeanExplore where available. If a physics API is missing, introduce faithful local abstractions rather than scalar placeholder aliases.

\begin{theorem}[Physics formalization target]
\label{thm:physics:phyx_mini_0944:target}
\lean{PhyXMiniProblems.ProblemPhyXMini0944.problem_phyx_mini_0944}
\uses{def:physics:phyx-mini-0944:phyxminiproblems-problemphyxmini0944-irradianceinwattspersquaremeter, def:physics:phyx-mini-0944:phyxminiproblems-problemphyxmini0944-electricamplitudeinvoltspermeter, def:physics:phyx-mini-0944:phyxminiproblems-problemphyxmini0944-magneticamplitudeinteslas, def:physics:phyx-mini-0944:phyxminiproblems-problemphyxmini0944-permeabilityinnewtonsperamperesquared, def:physics:phyx-mini-0944:phyxminiproblems-problemphyxmini0944-speedoflightinmeterspersecond, def:physics:phyx-mini-0944:phyxminiproblems-problemphyxmini0944-radiostationsatellitesetup, def:physics:phyx-mini-0944:phyxminiproblems-problemphyxmini0944-matchesproblemandprimaryfigure, def:physics:phyx-mini-0944:phyxminiproblems-problemphyxmini0944-hasphysicalradiationparameters, def:physics:phyx-mini-0944:phyxminiproblems-problemphyxmini0944-usestextbookvacuumpermeability, def:physics:phyx-mini-0944:phyxminiproblems-problemphyxmini0944-satisfiessinusoidaldetectionmodel, def:physics:phyx-mini-0944:phyxminiproblems-problemphyxmini0944-satisfieshemisphericalplanewavelaws, def:physics:phyx-mini-0944:phyxminiproblems-problemphyxmini0944-recordeddatasetanswer, def:physics:phyx-mini-0944:phyxminiproblems-problemphyxmini0944-magneticamplituderoundstochoiceb, def:physics:phyx-mini-0944:phyxminiproblems-problemphyxmini0944-electricamplituderoundstotwopointfourfivetimestenpownegtwo, def:physics:phyx-mini-0944:phyxminiproblems-problemphyxmini0944-isuniqueclosestmagneticamplitudechoice}
The assigned autoformalize agent should translate this physics problem into Lean declarations in the covered file, with theorem and lemma proof bodies written as `by sorry`.
\end{theorem}
\begin{proof}
This is an autoformalization task, not a proof task. Produce statements that can later be proved by the physics prover without weakening the physical model.
\end{proof}

% --- Archon declaration topology begin ---
\section{Declaration topology}
\begin{definition}[power Dimension]
  \label{def:physics:phyx-mini-0944:phyxminiproblems-problemphyxmini0944-powerdimension}
  \lean{PhyXMiniProblems.ProblemPhyXMini0944.powerDimension}
  The physical dimension M L² T⁻³ of power
\end{definition}
\begin{proof}
  This is the declared model definition of power Dimension.
\end{proof}
\begin{definition}[irradiance Dimension]
  \label{def:physics:phyx-mini-0944:phyxminiproblems-problemphyxmini0944-irradiancedimension}
  \lean{PhyXMiniProblems.ProblemPhyXMini0944.irradianceDimension}
  The physical dimension M T⁻³ of irradiance (power per area)
\end{definition}
\begin{proof}
  This is the declared model definition of irradiance Dimension.
\end{proof}
\begin{definition}[electric Field Strength Dimension]
  \label{def:physics:phyx-mini-0944:phyxminiproblems-problemphyxmini0944-electricfieldstrengthdimension}
  \lean{PhyXMiniProblems.ProblemPhyXMini0944.electricFieldStrengthDimension}
  The physical dimension M L T⁻² C⁻¹ of electric-field strength
\end{definition}
\begin{proof}
  This is the declared model definition of electric Field Strength Dimension.
\end{proof}
\begin{definition}[magnetic Field Strength Dimension]
  \label{def:physics:phyx-mini-0944:phyxminiproblems-problemphyxmini0944-magneticfieldstrengthdimension}
  \lean{PhyXMiniProblems.ProblemPhyXMini0944.magneticFieldStrengthDimension}
  The physical dimension M T⁻¹ C⁻¹ of magnetic-field strength
\end{definition}
\begin{proof}
  This is the declared model definition of magnetic Field Strength Dimension.
\end{proof}
\begin{definition}[magnetic Permeability Dimension]
  \label{def:physics:phyx-mini-0944:phyxminiproblems-problemphyxmini0944-magneticpermeabilitydimension}
  \lean{PhyXMiniProblems.ProblemPhyXMini0944.magneticPermeabilityDimension}
  The physical dimension M L C⁻² of vacuum magnetic permeability
\end{definition}
\begin{proof}
  This is the declared model definition of magnetic Permeability Dimension.
\end{proof}
\begin{definition}[Length Quantity]
  \label{def:physics:phyx-mini-0944:phyxminiproblems-problemphyxmini0944-lengthquantity}
  \lean{PhyXMiniProblems.ProblemPhyXMini0944.LengthQuantity}
  A nonnegative, unit-independent physical length
\end{definition}
\begin{proof}
  This is the declared model definition of Length Quantity.
\end{proof}
\begin{definition}[Power Quantity]
  \label{def:physics:phyx-mini-0944:phyxminiproblems-problemphyxmini0944-powerquantity}
  \lean{PhyXMiniProblems.ProblemPhyXMini0944.PowerQuantity}
  \uses{def:physics:phyx-mini-0944:phyxminiproblems-problemphyxmini0944-powerdimension}
  A nonnegative, unit-independent average power
\end{definition}
\begin{proof}
  This is the declared model definition of Power Quantity.
\end{proof}
\begin{definition}[Irradiance Quantity]
  \label{def:physics:phyx-mini-0944:phyxminiproblems-problemphyxmini0944-irradiancequantity}
  \lean{PhyXMiniProblems.ProblemPhyXMini0944.IrradianceQuantity}
  \uses{def:physics:phyx-mini-0944:phyxminiproblems-problemphyxmini0944-irradiancedimension}
  A nonnegative, unit-independent electromagnetic irradiance
\end{definition}
\begin{proof}
  This is the declared model definition of Irradiance Quantity.
\end{proof}
\begin{definition}[Electric Field Amplitude Quantity]
  \label{def:physics:phyx-mini-0944:phyxminiproblems-problemphyxmini0944-electricfieldamplitudequantity}
  \lean{PhyXMiniProblems.ProblemPhyXMini0944.ElectricFieldAmplitudeQuantity}
  \uses{def:physics:phyx-mini-0944:phyxminiproblems-problemphyxmini0944-electricfieldstrengthdimension}
  A nonnegative electric-field amplitude
\end{definition}
\begin{proof}
  This is the declared model definition of Electric Field Amplitude Quantity.
\end{proof}
\begin{definition}[Magnetic Field Amplitude Quantity]
  \label{def:physics:phyx-mini-0944:phyxminiproblems-problemphyxmini0944-magneticfieldamplitudequantity}
  \lean{PhyXMiniProblems.ProblemPhyXMini0944.MagneticFieldAmplitudeQuantity}
  \uses{def:physics:phyx-mini-0944:phyxminiproblems-problemphyxmini0944-magneticfieldstrengthdimension}
  A nonnegative magnetic-field amplitude
\end{definition}
\begin{proof}
  This is the declared model definition of Magnetic Field Amplitude Quantity.
\end{proof}
\begin{definition}[Magnetic Permeability Quantity]
  \label{def:physics:phyx-mini-0944:phyxminiproblems-problemphyxmini0944-magneticpermeabilityquantity}
  \lean{PhyXMiniProblems.ProblemPhyXMini0944.MagneticPermeabilityQuantity}
  \uses{def:physics:phyx-mini-0944:phyxminiproblems-problemphyxmini0944-magneticpermeabilitydimension}
  A nonnegative magnetic permeability, used here for vacuum μ₀
\end{definition}
\begin{proof}
  This is the declared model definition of Magnetic Permeability Quantity.
\end{proof}
\begin{definition}[length In Meters]
  \label{def:physics:phyx-mini-0944:phyxminiproblems-problemphyxmini0944-lengthinmeters}
  \lean{PhyXMiniProblems.ProblemPhyXMini0944.lengthInMeters}
  \uses{def:physics:phyx-mini-0944:phyxminiproblems-problemphyxmini0944-lengthquantity}
  Read a physical length in coherent SI metres
\end{definition}
\begin{proof}
  This is the declared model definition of length In Meters.
\end{proof}
\begin{definition}[length In Kilometers]
  \label{def:physics:phyx-mini-0944:phyxminiproblems-problemphyxmini0944-lengthinkilometers}
  \lean{PhyXMiniProblems.ProblemPhyXMini0944.lengthInKilometers}
  \uses{def:physics:phyx-mini-0944:phyxminiproblems-problemphyxmini0944-lengthquantity, def:physics:phyx-mini-0944:phyxminiproblems-problemphyxmini0944-lengthinmeters}
  Read a physical length in kilometres
\end{definition}
\begin{proof}
  This is the declared model definition of length In Kilometers.
\end{proof}
\begin{definition}[power In Watts]
  \label{def:physics:phyx-mini-0944:phyxminiproblems-problemphyxmini0944-powerinwatts}
  \lean{PhyXMiniProblems.ProblemPhyXMini0944.powerInWatts}
  \uses{def:physics:phyx-mini-0944:phyxminiproblems-problemphyxmini0944-powerquantity}
  Read an average power in coherent SI watts
\end{definition}
\begin{proof}
  This is the declared model definition of power In Watts.
\end{proof}
\begin{definition}[power In Kilowatts]
  \label{def:physics:phyx-mini-0944:phyxminiproblems-problemphyxmini0944-powerinkilowatts}
  \lean{PhyXMiniProblems.ProblemPhyXMini0944.powerInKilowatts}
  \uses{def:physics:phyx-mini-0944:phyxminiproblems-problemphyxmini0944-powerquantity, def:physics:phyx-mini-0944:phyxminiproblems-problemphyxmini0944-powerinwatts}
  Read an average power in kilowatts
\end{definition}
\begin{proof}
  This is the declared model definition of power In Kilowatts.
\end{proof}
\begin{definition}[irradiance In Watts Per Square Meter]
  \label{def:physics:phyx-mini-0944:phyxminiproblems-problemphyxmini0944-irradianceinwattspersquaremeter}
  \lean{PhyXMiniProblems.ProblemPhyXMini0944.irradianceInWattsPerSquareMeter}
  \uses{def:physics:phyx-mini-0944:phyxminiproblems-problemphyxmini0944-irradiancequantity}
  Read irradiance in coherent SI watts per square metre
\end{definition}
\begin{proof}
  This is the declared model definition of irradiance In Watts Per Square Meter.
\end{proof}
\begin{definition}[electric Amplitude In Volts Per Meter]
  \label{def:physics:phyx-mini-0944:phyxminiproblems-problemphyxmini0944-electricamplitudeinvoltspermeter}
  \lean{PhyXMiniProblems.ProblemPhyXMini0944.electricAmplitudeInVoltsPerMeter}
  \uses{def:physics:phyx-mini-0944:phyxminiproblems-problemphyxmini0944-electricfieldamplitudequantity}
  Read an electric-field amplitude in coherent SI volts per metre
\end{definition}
\begin{proof}
  This is the declared model definition of electric Amplitude In Volts Per Meter.
\end{proof}
\begin{definition}[magnetic Amplitude In Teslas]
  \label{def:physics:phyx-mini-0944:phyxminiproblems-problemphyxmini0944-magneticamplitudeinteslas}
  \lean{PhyXMiniProblems.ProblemPhyXMini0944.magneticAmplitudeInTeslas}
  \uses{def:physics:phyx-mini-0944:phyxminiproblems-problemphyxmini0944-magneticfieldamplitudequantity}
  Read a magnetic-field amplitude in coherent SI teslas
\end{definition}
\begin{proof}
  This is the declared model definition of magnetic Amplitude In Teslas.
\end{proof}
\begin{definition}[permeability In Newtons Per Ampere Squared]
  \label{def:physics:phyx-mini-0944:phyxminiproblems-problemphyxmini0944-permeabilityinnewtonsperamperesquared}
  \lean{PhyXMiniProblems.ProblemPhyXMini0944.permeabilityInNewtonsPerAmpereSquared}
  \uses{def:physics:phyx-mini-0944:phyxminiproblems-problemphyxmini0944-magneticpermeabilityquantity}
  Read vacuum permeability in coherent SI newtons per ampere squared
\end{definition}
\begin{proof}
  This is the declared model definition of permeability In Newtons Per Ampere Squared.
\end{proof}
\begin{definition}[speed Of Light In Meters Per Second]
  \label{def:physics:phyx-mini-0944:phyxminiproblems-problemphyxmini0944-speedoflightinmeterspersecond}
  \lean{PhyXMiniProblems.ProblemPhyXMini0944.speedOfLightInMetersPerSecond}
  Physlib's exact vacuum light speed, read in metres per second
\end{definition}
\begin{proof}
  This is the declared model definition of speed Of Light In Meters Per Second.
\end{proof}
\begin{definition}[Waveform]
  \label{def:physics:phyx-mini-0944:phyxminiproblems-problemphyxmini0944-waveform}
  \lean{PhyXMiniProblems.ProblemPhyXMini0944.Waveform}
  Qualitative waveform stated in the problem
\end{definition}
\begin{proof}
  This is the declared model definition of Waveform.
\end{proof}
\begin{definition}[Radiation Region]
  \label{def:physics:phyx-mini-0944:phyxminiproblems-problemphyxmini0944-radiationregion}
  \lean{PhyXMiniProblems.ProblemPhyXMini0944.RadiationRegion}
  The idealized angular region into which the transmitter radiates
\end{definition}
\begin{proof}
  This is the declared model definition of Radiation Region.
\end{proof}
\begin{definition}[Figure Object]
  \label{def:physics:phyx-mini-0944:phyxminiproblems-problemphyxmini0944-figureobject}
  \lean{PhyXMiniProblems.ProblemPhyXMini0944.FigureObject}
  The two labeled physical objects visible in the primary figure
\end{definition}
\begin{proof}
  This is the declared model definition of Figure Object.
\end{proof}
\begin{definition}[Radio Station Satellite Figure]
  \label{def:physics:phyx-mini-0944:phyxminiproblems-problemphyxmini0944-radiostationsatellitefigure}
  \lean{PhyXMiniProblems.ProblemPhyXMini0944.RadioStationSatelliteFigure}
  \uses{def:physics:phyx-mini-0944:phyxminiproblems-problemphyxmini0944-lengthquantity, def:physics:phyx-mini-0944:phyxminiproblems-problemphyxmini0944-figureobject}
  Literal typed content of image 944.png. The radius label is a physical length. Its 100 km calibration is imposed separately, rather than stored as a dimensionless field of the figure
\end{definition}
\begin{proof}
  This is the declared model definition of Radio Station Satellite Figure.
\end{proof}
\begin{definition}[Radio Station Satellite Setup]
  \label{def:physics:phyx-mini-0944:phyxminiproblems-problemphyxmini0944-radiostationsatellitesetup}
  \lean{PhyXMiniProblems.ProblemPhyXMini0944.RadioStationSatelliteSetup}
  \uses{def:physics:phyx-mini-0944:phyxminiproblems-problemphyxmini0944-lengthquantity, def:physics:phyx-mini-0944:phyxminiproblems-problemphyxmini0944-powerquantity, def:physics:phyx-mini-0944:phyxminiproblems-problemphyxmini0944-irradiancequantity, def:physics:phyx-mini-0944:phyxminiproblems-problemphyxmini0944-electricfieldamplitudequantity, def:physics:phyx-mini-0944:phyxminiproblems-problemphyxmini0944-magneticfieldamplitudequantity, def:physics:phyx-mini-0944:phyxminiproblems-problemphyxmini0944-magneticpermeabilityquantity, def:physics:phyx-mini-0944:phyxminiproblems-problemphyxmini0944-waveform, def:physics:phyx-mini-0944:phyxminiproblems-problemphyxmini0944-radiationregion, def:physics:phyx-mini-0944:phyxminiproblems-problemphyxmini0944-radiostationsatellitefigure}
  The independent source, detector, fields, and amplitude observables. Neither amplitude is defined from a displayed answer or from the requested numerical value
\end{definition}
\begin{proof}
  This is the declared model definition of Radio Station Satellite Setup.
\end{proof}
\begin{definition}[expected Figure Label]
  \label{def:physics:phyx-mini-0944:phyxminiproblems-problemphyxmini0944-expectedfigurelabel}
  \lean{PhyXMiniProblems.ProblemPhyXMini0944.expectedFigureLabel}
  \uses{def:physics:phyx-mini-0944:phyxminiproblems-problemphyxmini0944-figureobject}
  The literal labels printed next to each object in the supplied image
\end{definition}
\begin{proof}
  This is the declared model definition of expected Figure Label.
\end{proof}
\begin{definition}[Matches Problem And Primary Figure]
  \label{def:physics:phyx-mini-0944:phyxminiproblems-problemphyxmini0944-matchesproblemandprimaryfigure}
  \lean{PhyXMiniProblems.ProblemPhyXMini0944.MatchesProblemAndPrimaryFigure}
  \uses{def:physics:phyx-mini-0944:phyxminiproblems-problemphyxmini0944-lengthinmeters, def:physics:phyx-mini-0944:phyxminiproblems-problemphyxmini0944-lengthinkilometers, def:physics:phyx-mini-0944:phyxminiproblems-problemphyxmini0944-powerinkilowatts, def:physics:phyx-mini-0944:phyxminiproblems-problemphyxmini0944-radiostationsatellitesetup, def:physics:phyx-mini-0944:phyxminiproblems-problemphyxmini0944-expectedfigurelabel}
  Problem-text and primary-image readouts. These facts identify the source power and range and relate the physical setup to the dome, labels, and arrow. They do not assign either field amplitude
\end{definition}
\begin{proof}
  This is the declared model definition of Matches Problem And Primary Figure.
\end{proof}
\begin{definition}[Has Physical Radiation Parameters]
  \label{def:physics:phyx-mini-0944:phyxminiproblems-problemphyxmini0944-hasphysicalradiationparameters}
  \lean{PhyXMiniProblems.ProblemPhyXMini0944.HasPhysicalRadiationParameters}
  \uses{def:physics:phyx-mini-0944:phyxminiproblems-problemphyxmini0944-lengthinmeters, def:physics:phyx-mini-0944:phyxminiproblems-problemphyxmini0944-powerinwatts, def:physics:phyx-mini-0944:phyxminiproblems-problemphyxmini0944-irradianceinwattspersquaremeter, def:physics:phyx-mini-0944:phyxminiproblems-problemphyxmini0944-permeabilityinnewtonsperamperesquared, def:physics:phyx-mini-0944:phyxminiproblems-problemphyxmini0944-speedoflightinmeterspersecond, def:physics:phyx-mini-0944:phyxminiproblems-problemphyxmini0944-radiostationsatellitesetup}
  Positivity and nondegeneracy conditions for the radiation model
\end{definition}
\begin{proof}
  This is the declared model definition of Has Physical Radiation Parameters.
\end{proof}
\begin{definition}[Uses Textbook Vacuum Permeability]
  \label{def:physics:phyx-mini-0944:phyxminiproblems-problemphyxmini0944-usestextbookvacuumpermeability}
  \lean{PhyXMiniProblems.ProblemPhyXMini0944.UsesTextbookVacuumPermeability}
  \uses{def:physics:phyx-mini-0944:phyxminiproblems-problemphyxmini0944-permeabilityinnewtonsperamperesquared, def:physics:phyx-mini-0944:phyxminiproblems-problemphyxmini0944-radiostationsatellitesetup}
  The standard textbook calibration μ₀ = 4π × 10⁻⁷ N/A². This is an independent universal-constant input, not a requested field amplitude
\end{definition}
\begin{proof}
  This is the declared model definition of Uses Textbook Vacuum Permeability.
\end{proof}
\begin{definition}[Satisfies Sinusoidal Detection Model]
  \label{def:physics:phyx-mini-0944:phyxminiproblems-problemphyxmini0944-satisfiessinusoidaldetectionmodel}
  \lean{PhyXMiniProblems.ProblemPhyXMini0944.SatisfiesSinusoidalDetectionModel}
  \uses{def:physics:phyx-mini-0944:phyxminiproblems-problemphyxmini0944-electricamplitudeinvoltspermeter, def:physics:phyx-mini-0944:phyxminiproblems-problemphyxmini0944-magneticamplitudeinteslas, def:physics:phyx-mini-0944:phyxminiproblems-problemphyxmini0944-radiostationsatellitesetup}
  At the satellite, both fields are sinusoidal with a common phase and with the stored amplitudes as their SI component scales. Unit polarization vectors ensure that these scales really are field maxima; no numerical amplitude is assumed
\end{definition}
\begin{proof}
  This is the declared model definition of Satisfies Sinusoidal Detection Model.
\end{proof}
\begin{definition}[Satisfies Hemispherical Plane Wave Laws]
  \label{def:physics:phyx-mini-0944:phyxminiproblems-problemphyxmini0944-satisfieshemisphericalplanewavelaws}
  \lean{PhyXMiniProblems.ProblemPhyXMini0944.SatisfiesHemisphericalPlaneWaveLaws}
  \uses{def:physics:phyx-mini-0944:phyxminiproblems-problemphyxmini0944-lengthinmeters, def:physics:phyx-mini-0944:phyxminiproblems-problemphyxmini0944-powerinwatts, def:physics:phyx-mini-0944:phyxminiproblems-problemphyxmini0944-irradianceinwattspersquaremeter, def:physics:phyx-mini-0944:phyxminiproblems-problemphyxmini0944-electricamplitudeinvoltspermeter, def:physics:phyx-mini-0944:phyxminiproblems-problemphyxmini0944-magneticamplitudeinteslas, def:physics:phyx-mini-0944:phyxminiproblems-problemphyxmini0944-permeabilityinnewtonsperamperesquared, def:physics:phyx-mini-0944:phyxminiproblems-problemphyxmini0944-speedoflightinmeterspersecond, def:physics:phyx-mini-0944:phyxminiproblems-problemphyxmini0944-radiostationsatellitesetup}
  The equal-above-ground radiation assumption spreads the average power over a hemisphere of area 2πr². A vacuum plane wave has E\_max = c B\_max, and its time-averaged Poynting magnitude is E\_max B\_max / (2 μ₀). These are general modeling relations and contain none of the requested numerical conclusions
\end{definition}
\begin{proof}
  This is the declared model definition of Satisfies Hemispherical Plane Wave Laws.
\end{proof}
\begin{lemma}[magnetic Amplitude exact]
  \label{lem:physics:phyx-mini-0944:phyxminiproblems-problemphyxmini0944-magneticamplitude-exact}
  \lean{PhyXMiniProblems.ProblemPhyXMini0944.magneticAmplitude_exact}
  \uses{def:physics:phyx-mini-0944:phyxminiproblems-problemphyxmini0944-irradianceinwattspersquaremeter, def:physics:phyx-mini-0944:phyxminiproblems-problemphyxmini0944-magneticamplitudeinteslas, def:physics:phyx-mini-0944:phyxminiproblems-problemphyxmini0944-permeabilityinnewtonsperamperesquared, def:physics:phyx-mini-0944:phyxminiproblems-problemphyxmini0944-speedoflightinmeterspersecond, def:physics:phyx-mini-0944:phyxminiproblems-problemphyxmini0944-radiostationsatellitesetup, def:physics:phyx-mini-0944:phyxminiproblems-problemphyxmini0944-hasphysicalradiationparameters, def:physics:phyx-mini-0944:phyxminiproblems-problemphyxmini0944-satisfieshemisphericalplanewavelaws}
  Solving the positive plane-wave and Poynting relations for B\_max gives the standard exact square-root expression
\end{lemma}
\begin{proof}
  The conclusion follows from the governing relations and figure readouts named in the statement of magnetic Amplitude exact.
\end{proof}
\begin{lemma}[electric Amplitude exact]
  \label{lem:physics:phyx-mini-0944:phyxminiproblems-problemphyxmini0944-electricamplitude-exact}
  \lean{PhyXMiniProblems.ProblemPhyXMini0944.electricAmplitude_exact}
  \uses{def:physics:phyx-mini-0944:phyxminiproblems-problemphyxmini0944-irradianceinwattspersquaremeter, def:physics:phyx-mini-0944:phyxminiproblems-problemphyxmini0944-electricamplitudeinvoltspermeter, def:physics:phyx-mini-0944:phyxminiproblems-problemphyxmini0944-permeabilityinnewtonsperamperesquared, def:physics:phyx-mini-0944:phyxminiproblems-problemphyxmini0944-speedoflightinmeterspersecond, def:physics:phyx-mini-0944:phyxminiproblems-problemphyxmini0944-radiostationsatellitesetup, def:physics:phyx-mini-0944:phyxminiproblems-problemphyxmini0944-hasphysicalradiationparameters, def:physics:phyx-mini-0944:phyxminiproblems-problemphyxmini0944-satisfieshemisphericalplanewavelaws}
  The corresponding exact vacuum relation for E\_max
\end{lemma}
\begin{proof}
  The conclusion follows from the governing relations and figure readouts named in the statement of electric Amplitude exact.
\end{proof}
\begin{definition}[Answer Choice]
  \label{def:physics:phyx-mini-0944:phyxminiproblems-problemphyxmini0944-answerchoice}
  \lean{PhyXMiniProblems.ProblemPhyXMini0944.AnswerChoice}
  Labels of the four magnetic-field choices printed in the question
\end{definition}
\begin{proof}
  This is the declared model definition of Answer Choice.
\end{proof}
\begin{definition}[displayed Magnetic Amplitude In Teslas]
  \label{def:physics:phyx-mini-0944:phyxminiproblems-problemphyxmini0944-displayedmagneticamplitudeinteslas}
  \lean{PhyXMiniProblems.ProblemPhyXMini0944.displayedMagneticAmplitudeInTeslas}
  \uses{def:physics:phyx-mini-0944:phyxminiproblems-problemphyxmini0944-answerchoice}
  Magnetic-field amplitude in teslas printed beside each answer label
\end{definition}
\begin{proof}
  This is the declared model definition of displayed Magnetic Amplitude In Teslas.
\end{proof}
\begin{definition}[recorded Dataset Answer]
  \label{def:physics:phyx-mini-0944:phyxminiproblems-problemphyxmini0944-recordeddatasetanswer}
  \lean{PhyXMiniProblems.ProblemPhyXMini0944.recordedDatasetAnswer}
  \uses{def:physics:phyx-mini-0944:phyxminiproblems-problemphyxmini0944-answerchoice}
  The answer label recorded in the source dataset
\end{definition}
\begin{proof}
  This is the declared model definition of recorded Dataset Answer.
\end{proof}
\begin{definition}[magnetic Choice BTolerance In Teslas]
  \label{def:physics:phyx-mini-0944:phyxminiproblems-problemphyxmini0944-magneticchoicebtoleranceinteslas}
  \lean{PhyXMiniProblems.ProblemPhyXMini0944.magneticChoiceBToleranceInTeslas}
  Half of the final displayed digit in the three-digit value of choice B
\end{definition}
\begin{proof}
  This is the declared model definition of magnetic Choice BTolerance In Teslas.
\end{proof}
\begin{definition}[electric Amplitude Tolerance In Volts Per Meter]
  \label{def:physics:phyx-mini-0944:phyxminiproblems-problemphyxmini0944-electricamplitudetoleranceinvoltspermeter}
  \lean{PhyXMiniProblems.ProblemPhyXMini0944.electricAmplitudeToleranceInVoltsPerMeter}
  Half of the final displayed digit in 2.45 × 10⁻² V/m
\end{definition}
\begin{proof}
  This is the declared model definition of electric Amplitude Tolerance In Volts Per Meter.
\end{proof}
\begin{definition}[Magnetic Amplitude Rounds To Choice B]
  \label{def:physics:phyx-mini-0944:phyxminiproblems-problemphyxmini0944-magneticamplituderoundstochoiceb}
  \lean{PhyXMiniProblems.ProblemPhyXMini0944.MagneticAmplitudeRoundsToChoiceB}
  \uses{def:physics:phyx-mini-0944:phyxminiproblems-problemphyxmini0944-magneticamplitudeinteslas, def:physics:phyx-mini-0944:phyxminiproblems-problemphyxmini0944-radiostationsatellitesetup, def:physics:phyx-mini-0944:phyxminiproblems-problemphyxmini0944-displayedmagneticamplitudeinteslas, def:physics:phyx-mini-0944:phyxminiproblems-problemphyxmini0944-magneticchoicebtoleranceinteslas}
  A magnetic amplitude rounds to the displayed value of choice B
\end{definition}
\begin{proof}
  This is the declared model definition of Magnetic Amplitude Rounds To Choice B.
\end{proof}
\begin{definition}[Electric Amplitude Rounds To Two Point Four Five Times Ten Pow Neg Two]
  \label{def:physics:phyx-mini-0944:phyxminiproblems-problemphyxmini0944-electricamplituderoundstotwopointfourfivetimestenpownegtwo}
  \lean{PhyXMiniProblems.ProblemPhyXMini0944.ElectricAmplitudeRoundsToTwoPointFourFiveTimesTenPowNegTwo}
  \uses{def:physics:phyx-mini-0944:phyxminiproblems-problemphyxmini0944-electricamplitudeinvoltspermeter, def:physics:phyx-mini-0944:phyxminiproblems-problemphyxmini0944-radiostationsatellitesetup, def:physics:phyx-mini-0944:phyxminiproblems-problemphyxmini0944-electricamplitudetoleranceinvoltspermeter}
  An electric amplitude rounds to 2.45 × 10⁻² V/m
\end{definition}
\begin{proof}
  This is the declared model definition of Electric Amplitude Rounds To Two Point Four Five Times Ten Pow Neg Two.
\end{proof}
\begin{definition}[Is Unique Closest Magnetic Amplitude Choice]
  \label{def:physics:phyx-mini-0944:phyxminiproblems-problemphyxmini0944-isuniqueclosestmagneticamplitudechoice}
  \lean{PhyXMiniProblems.ProblemPhyXMini0944.IsUniqueClosestMagneticAmplitudeChoice}
  \uses{def:physics:phyx-mini-0944:phyxminiproblems-problemphyxmini0944-magneticamplitudeinteslas, def:physics:phyx-mini-0944:phyxminiproblems-problemphyxmini0944-radiostationsatellitesetup, def:physics:phyx-mini-0944:phyxminiproblems-problemphyxmini0944-answerchoice, def:physics:phyx-mini-0944:phyxminiproblems-problemphyxmini0944-displayedmagneticamplitudeinteslas}
  A choice is strictly closer than every distinct displayed alternative
\end{definition}
\begin{proof}
  This is the declared model definition of Is Unique Closest Magnetic Amplitude Choice.
\end{proof}
% --- Archon declaration topology end ---
% --- Archon physics formalization source end ---
